\pdfminorversion=7%
\documentclass[aspectratio=169,mathserif,notheorems]{beamer}%
%
\xdef\bookbaseDir{../../bookbase}%
\xdef\sharedDir{../../shared}%
\RequirePackage{\bookbaseDir/styles/slides}%
\RequirePackage{\sharedDir/styles/styles}%
\toggleToGerman%
%
\subtitle{7.~Beispiele Herunterladen}%
%
\begin{document}%
%
\startPresentation%
%
\section{Einleitung}%
%
\begin{frame}%
\frametitle{Einleitung}%
\begin{itemize}%
\item Im Rest dieses Kurses werden wir intensiv mit praktischen Beispielen arbeiten.%
\item<2-> Wenn wir ein Konzept vorstellen, dann probieren wir das praktisch aus.%
\item<3-> Damit Sie diese Beispiele nachvollziehen können, stellen wir sie in einem \github\ repository zur Verfügung.%
\item<4-> Dieses Repository finden Sie unter \expandafter\url{\databasesCodeRepo}.%
\end{itemize}%
\end{frame}%
%
\section{Herunterladen der Beispiele}%
%
\begin{frame}[t]%
\frametitle{Herunterladen der Beispiele}%
\begin{itemize}%
\only<-2>{\item Öffnen Sie einen Webbrowser und besuchen Sie die Webseite \expandafter\url{\databasesCodeRepo}.%
\only<2>{ Klicken Sie auf das nach unten gerichtete Dreieck im Button \menu{Code}.}%
}%
%
\only<3>{\item<3-> In dem sich öffnenden Menü, klicken Sie auf \menu{Download ZIP}.}%
%
\only<4-5>{\item<4-> Der Download beginnt und ist irgendwann abgeschlossen.%
\only<5->{ Öffnen Sie die heruntegrladene Datei.}}%
\only<6->{\item<6-> Die Datei ist ein ZIP-Archiv, also eine Datei, die andere Dateien und Verzeichnisse enthält.%
\only<7->{ Sie können Sie an einem Ihnen angenehmen Ort entpacken. %
\only<8->{ Das erste große Beispiel, mit dem wir als nächstes beginnen, befindet sich üprigens im Unterordner \textil{factory}.}%
}%
}%
\end{itemize}%
%
\locateGraphicTB{-2}{width=0.6\paperwidth}{graphics/examplesBrowser1website}{0.2}{0.33}%
\locateGraphicTB{3}{width=0.6\paperwidth}{graphics/examplesBrowser2beginDownload}{0.2}{0.33}%
\locateGraphicTB{4-5}{width=0.6\paperwidth}{graphics/examplesBrowser3downloaded}{0.2}{0.33}%
\locateGraphicTB{6-}{width=0.6\paperwidth}{graphics/examplesBrowser4opened}{0.2}{0.33}%
\end{frame}%
%
\section{Zusammenfassung}%
%
\begin{frame}%
\frametitle{Zusammenfassung}%
\begin{itemize}%
\item Wir haben nun das Repository mit den Beispielen für diesen Kurs heruntergeladen.%
\item<2-> Damit haben Sie alle Skripte, die wir im folgenden verwenden, direkt zur Hand.%
\item<3-> Sie können also unsere Beispiele sehr komfortabel nachvollziehen.%
\end{itemize}%
\end{frame}%
%
\endPresentation%
\end{document}%%
\endinput%
%
